\section*{Bab \ref{ch:b2:genome-diversification} --- Mutasi dan Penganekaragaman Genom}

\begin{solution}{exo:b2:genome-diversification:1}
GAG: bisu (tetap glutamat), sebuah transisi (A$\to$G). GTA: salah arti
(valina), sebuah transversi (A$\to$T). TAA: tanpa arti (henti), sebuah
transversi (G$\to$T). Menyisipkan sebuah basa: sebuah \omterm{def:b2:genome-diversification:mutation}{pergeseran rangka
baca}, yang menulis ulang setiap kodon sesudahnya.
\end{solution}

\begin{solution}{exo:b2:genome-diversification:2}
$U = 5\times 10^{-10}\times 4.6\times 10^{6} = 2.3\times 10^{-3}$: satu
pembelahan dari kira-kira 430 menghasilkan sebuah \omterm{def:b2:genome-diversification:mutation}{mutasi}. Dalam
$10^{9}$ sel: $2.3\times 10^{6}$ \omterm{def:b2:genome-diversification:mutation}{mutasi} baru tiap generasi.
\end{solution}

\begin{solution}{exo:b2:genome-diversification:3}
\omterm{def:b2:genome-diversification:hgt}{Transformasi}: penyerapan DNA bebas dari sel yang lisis, membutuhkan
seorang penerima yang berkompeten; memindahkan serpihan kromosom (gen
kapsul pada pneumokokus). Konjugasi: sentuhan antarsel lewat sebuah
pilus, membutuhkan sebuah \omterm{def:b2:unicellular-diversity:prokaryote}{plasmid} konjugatif pada pemberinya;
memindahkan \omterm{def:b2:unicellular-diversity:prokaryote}{plasmid} (gen kekebalan) dan, dari sebuah Hfr, gen
kromosom. \omterm{def:b2:genome-diversification:hgt}{Transduksi}: sebuah bakteriofag yang mengemas DNA inangnya;
memindahkan serpihan acak (terampat) atau gen di sebelah tapak
profagnya (terkhusus).
\end{solution}

\begin{solution}{exo:b2:genome-diversification:4}
\omterm{def:b2:genome-diversification:polyploidy}{Autopoliploid}: perangkat tambahan dari satu spesies, lewat sebuah gamet
diploid (kentang tetraploid, banyak semanggi autotetraploid).
\omterm{def:b2:genome-diversification:polyploidy}{Alopoliploid}: perangkat dua spesies yang digabungkan lewat penghibridan
lalu digandakan (gandum roti, kapas, tembakau). Hibrid diploid AB tidak
memiliki dua kromosom yang serupa, sehingga tak satu pun dapat
berpasangan pada meiosis I; kromosomnya memilah secara acak dan
gametnya tak berimbang serta tak sanggup hidup.
\end{solution}

\begin{solution}{exo:b2:genome-diversification:5}
$2\times 3.2\times 10^{9}\times 1.2\times 10^{-8} = 77$ \omterm{def:b2:genome-diversification:mutation}{mutasi} baru.
Di dalam urutan penyandi: $77\times 0.015 = 1.2$; yang mengubah sebuah
asam amino: $1.2\times 0.7 = 0.8$ --- sebagian besar anak mengusung
kira-kira satu substitusi asam amino yang baru.
\end{solution}

\begin{solution}{exo:b2:genome-diversification:6}
Peristiwanya: $m(N - N_0) = 2\times 10^{-8}\times 10^{9} = 20$. Sel
mutannya: $mN\ln(N/N_0) = 20\times \ln 10^{6} = 20\times 13.8 = 276$.
$p_0 = e^{-20} \approx 2\times 10^{-9}$: tiap biakannya bermutan,
sehingga mencacah biakan yang tak bermutan tidak mengukur apa pun. Yang
dibutuhkan adalah $m(N -
N_0) \approx 1$, yakni $N \approx 5\times 10^{7}$, tempat $p_0 \approx
0.37$ dan beberapa lusin biakan memberikan $m$ dengan ketelitian sampai
faktor dua.
\end{solution}

\begin{solution}{exo:b2:genome-diversification:7}
Urasil tidak termasuk ke dalam DNA, sehingga sebuah urasil-DNA
glikosilase dapat membuang setiap urasil yang ditemukannya tanpa
kesamaran. Timina adalah basa normal: sesudah 5-metilsitosina
terdeaminasi, salah pasang G$\cdot$T memang terlihat, tetapi mesin
perbaikannya tidak punya cara mengetahui untai mana yang keliru dan
separuh waktunya ``membetulkan'' G-nya, sehingga menetapkan sebuah
transisi C$\to$T. Maka tapak CpG yang termetilasi bermutasi kira-kira
sepuluh kali lebih cepat daripada tapak lain, dan sepanjang waktu
evolusi mereka telah berubah menjadi TpG dan CpA: genom mamalia hanya
memuat kira-kira seperlima CpG yang diramalkan susunan basanya.
\end{solution}

\begin{solution}{exo:b2:genome-diversification:8}
\qty{46}{kb} tiap menit. Kedudukannya: 368, 782, 1150, dan
\qty{1840}{kb}; jaraknya 414, 368, dan \qty{690}{kb}.
\end{solution}

\begin{solution}{exo:b2:genome-diversification:9}
Salinan kedua satu kromosom berpasangan dengan salinan pertama kromosom
yang lain; sebuah pindah silang di antaranya memberi satu kromatid yang
mengusung salinan~1, salinan~2, dan salinan~2$'$ (tiga gen) dan satu
kromatid yang hanya mengusung salinan~1$'$. Seseorang dengan kromosom
bergen satu dari tiap tetuanya memiliki dua gen $\alpha$ alih-alih
empat ($-\alpha/-\alpha$): rantai $\alpha$-nya dibuat pada laju
separuhnya, sel darah merahnya kecil dan pucat, dengan anemia ringan
atau tanpa anemia --- sifat pembawa talasemia $\alpha$.
\end{solution}

\begin{solution}{exo:b2:genome-diversification:10}
Gametnya: $n = 21$ (satu perangkat masing-masing dari A, B, dan D).
Hibridnya: AB, 14 kromosom; ABD, 21 kromosom. Gandum roti $\times$
gandum emer: AABBD, 35 kromosom --- perangkat A dan B berpasangan,
perangkat D tidak memiliki pasangan, dan tumbuhannya sebagian besar
mandul. Tiap penggandaan memberi tiap kromosomnya sebuah homolog yang
serupa, sehingga bivalennya terbentuk dan gametnya berimbang;
menyilangbalikkannya kepada seorang tetua menghasilkan tumbuhan dengan
satu perangkat tak berpasangan, yang gametnya tak berimbang, sehingga
si heksaploid terpencil secara perkembangbiakan dari moyangnya ---
sebuah spesies dalam satu langkah.
\end{solution}

\begin{solution}{exo:b2:genome-diversification:11}
Pertumbuhan logistik: $R(t) = N/\bigl(1 + (N - 1)e^{-\gamma Nt}\bigr)$.
Separuh populasinya ketika $(N - 1)e^{-\gamma Nt} = 1$: $t = \ln(N -
1)/\gamma N = 20.7/0.5 = \qty{41}{h}$. Tanpa antibiotiknya, sel yang
bebas \omterm{def:b2:unicellular-diversity:prokaryote}{plasmid} tumbuh \qty{5}{\%} lebih cepat, sehingga rasio pengusung
terhadap yang bukan pengusung meluruh seperti $e^{-st}$ dengan
$s = 0.05\ln 2 =
\qty{0.035}{h^{-1}}$ tiap jam (\omterm{def:b2:unicellular-diversity:division}{waktu generasi}
\qty{1}{h}): dari \qty{99}{\%} menjadi \qty{1}{\%} memakan
$\ln(99^2)/0.035 = \qty{260}{h}$, sebelas hari --- dan pemindahan yang
berlanjut dapat mempertahankan \omterm{def:b2:unicellular-diversity:prokaryote}{plasmidnya} tanpa batas pada frekuensi
rendah, siap bagi babak pengobatan berikutnya.
\end{solution}

\begin{solution}{exo:b2:genome-diversification:12}
\omterm{thm:b2:genome-diversification:luria}{Uji fluktuasi} menunjukkan bahwa \omterm{def:b2:genome-diversification:mutation}{mutasi} kekebalan muncul sebelum dan
tanpa zat penyeleksinya, pada waktu yang acak, sehingga \omterm{def:b2:genome-diversification:mutation}{mutasi} buta
terhadap apa yang akan berguna; pelempengan replika menunjukkan hal
yang sama tanpa memaparkan sel terseleksinya sama sekali. Evolusi tidak
acak karena seleksi memilah varian buta ini menurut akibatnya, secara
ajek dan tertimbun. Satu-satunya tempat keacakan itu sendiri dibentuk
adalah lajunya: garis keturunan dengan kekeliruan yang terlalu banyak
atau terlalu sedikit sama-sama kalah, sehingga \omterm{thm:b2:genome-diversification:rate}{laju mutasinya} adalah
kompromi yang teradaptasi --- tetapi apa yang dikerjakan tiap \omterm{def:b2:genome-diversification:mutation}{mutasi}
tetap tak teramalkan dan tak terpilihkan.
\end{solution}

\begin{solution}{pb:b2:genome-diversification:1}
\textbf{1.} Jumlahnya 119, puratanya $5.95$; purata kuadratnya
$11489/20 = 574$; ragamnya $574 - 35 = 540$.
\textbf{2.} Jumlahnya 324, puratanya $16.2$; purata kuadratnya
$5462/20 = 273$; ragamnya $273 - 262 = 11$.
\textbf{3.} Contoh dari satu biakan bersifat Poisson (ragamnya dekat
dengan puratanya): mereka adalah cabutan acak dari satu bagian mutan
yang tetap. Biakan yang saling bebas memiliki ragam sembilan puluh kali
puratanya: mutannya muncul pada waktu yang berbeda-beda dan acak selama
pertumbuhannya, dan makin dini \omterm{def:b2:genome-diversification:mutation}{mutasinya} makin besar klonnya.
\textbf{4.} Sebuah jekpot: sebuah \omterm{def:b2:genome-diversification:mutation}{mutasi} pada salah satu dari beberapa
pembelahan pertamanya, yang klonnya lalu tumbuh bersama biakannya
sampai seratus sel.
\textbf{5.} Empat belas dari dua puluh: $p_0 = 0.70$; $m(N - N_0) = -\ln 0.70
= 0.36$; $m = 0.36/2\times 10^{8} = \qty{1.8e-9}{}$ tiap pembelahan.
\textbf{6.} $1.8\times 10^{-9}\times 2\times 10^{8}\times \ln(2\times
10^{6}) = 0.36\times 14.5 = 5.2$ sel mutan; amatannya $5.95$ ---
sepadan dalam derau dua puluh biakan.
\textbf{7.} Tidak ada contoh yang kosong (semuanya berbagi mutan biakan
induknya), sehingga $p_0 = 0$ tidak memberi taksiran; contohnya
mengukur bagian mutan satu biakan, yang bergantung pada kapan \omterm{def:b2:genome-diversification:mutation}{mutasinya}
kebetulan terjadi.
\textbf{8.} $u = m/20 = \qty{9e-11}{}$ tiap pasangan basa tiap replikasi.
\textbf{9.} $U = 9\times 10^{-11}\times 4.6\times 10^{6} =
\qty{4.1e-4}{}$.
\textbf{10.} $10^{12}$ sel $\times$ $4.1\times 10^{-4}$ $= 4.1\times
10^{8}$ \omterm{def:b2:genome-diversification:mutation}{mutasi} baru.
\textbf{11.} $3\times 4.6\times 10^{6} = 1.4\times 10^{7}$ perubahan
yang mungkin; masing-masing muncul
$4.1\times 10^{8}/1.4\times 10^{7} \approx 30$ kali.
\textbf{12.} Setiap perubahan basa tunggal, termasuk setiap perubahan
yang memberikan kekebalan terhadap obat yang dapat dikalahkan satu
\omterm{def:b2:genome-diversification:mutation}{mutasi}, hadir dalam puluhan sel sebelum obatnya ditambahkan; obatnya
sekadar menyeleksi mereka.
\textbf{13.} Dua \omterm{def:b2:genome-diversification:mutation}{mutasi} yang saling bebas di dalam sel yang sama
terjadi pada $m^2 \approx (2\times 10^{-9})^2 = 4\times 10^{-18}$ tiap
pembelahan: dalam $10^{12}$ sel, $4\times 10^{-6}$ tiap generasi ---
sel yang kebal ganda pada dasarnya tidak pernah ada sebelumnya, dan
itulah sebabnya tuberkulosis diobati dengan beberapa obat sekaligus.
\textbf{14.} Dengan $a = N - 1$ dan $E = e^{-\gamma Nt}$: $R = N/(1 +
aE)$, $\mathrm{d}R/\mathrm{d}t = \gamma N\cdot NaE/(1 + aE)^2$; dan
$\gamma R(N - R) = \gamma\,\frac{N}{1 + aE}\cdot\frac{NaE}{1 + aE}$,
yang sama; $R(0) = N/(1 + N - 1) = 1$.
\textbf{15.} $(N - 1)E = 1$: $t = \ln(N - 1)/\gamma N = 20.7/0.5 =
\qty{41}{h}$.
\textbf{16.} $1 + aE = 1/0.99$, $aE = 0.0101$: $t = \ln(10^{9}/0.0101)
/0.5 = (20.7 + 4.6)/0.5 = \qty{51}{h}$.
\textbf{17.} $\gamma (N/2)^2 = \gamma N\cdot N/4 = 0.5\times 2.5\times
10^{8} = 1.25\times 10^{8}$ pemindahan tiap jam.
\textbf{18.} Laju pertumbuhannya $\ln 2 = \qty{0.693}{h^{-1}}$ dan
\qty{0.762}{h^{-1}}; rasio mutan terhadap yang peka tumbuh seperti
$e^{st}$ dengan $s = \qty{0.069}{h^{-1}}$; dari $1/N$ menjadi 1
memakan $\ln N/s =
20.7/0.069 = \qty{300}{h}$, dua belas hari.
\textbf{19.} Pemindahannya sebanding dengan hasil kali pengusung dan
yang bukan pengusung serta berjalan pada laju perjumpaan, bukan pada
laju pembelahan; pewarisan dibatasi keunggulan pertumbuhan mutannya,
yang kecil. Sebuah \omterm{def:b2:unicellular-diversity:prokaryote}{plasmid} melintasi populasinya dalam dua hari,
seorang mutan yang diunggulkan dalam dua pekan.
\textbf{20.} $U = 4.1\times 10^{-2}$; yang merugikan: $0.041\times 0.88\times
0.3 = 0.011$ tiap generasi.
\textbf{21.} Puratanya $2.2$: $e^{-2.2} = 0.11$, satu keturunan dari
sembilan selamat tanpa cedera.
\textbf{22.} Tipe liarnya: $1.1\times 10^{-4}$ tiap generasi, $0.022$
dalam 200 generasi: $e^{-0.022} = 0.98$.
\textbf{23.} Puratanya $520$ sel mutan; $p_0 = e^{-36} \approx
10^{-16}$: tidak ada biakan yang kosong.
\textbf{24.} Di bawah seleksi yang kuat dan berubah-ubah (antibiotik
rumah sakit yang bergilir), pasokan varian baru yang seratus kali lipat
mengalahkan bebannya, dan mutatornya menumpang bersama kekebalan yang
dihasilkannya. Di lingkungan yang mantap tidak ada yang dapat diraih
dan ada satu \omterm{def:b2:genome-diversification:mutation}{mutasi} merugikan tiap seratus generasi yang dapat
hilang, sehingga garis keturunan mutator dibersihkan.
\textbf{25.} $m = \qty{1.8e-9}{}$ tiap pembelahan menuju kekebalan, $u =
\qty{9e-11}{}$ tiap pasangan basa tiap replikasi,
$U = \qty{4.1e-4}{}$ \omterm{def:b2:genome-diversification:mutation}{mutasi} tiap genom tiap replikasi; \omterm{def:b2:unicellular-diversity:prokaryote}{plasmidnya}
mencapai separuh ususnya dalam \qty{41}{h} dan \qty{99}{\%} ususnya
dalam \qty{51}{h}.
\end{solution}
